\section{Introduction}

MASFE (Multi-Algorithm Scheduling and Fusion Engine) answers NASA's Space-to-Soil challenge by moving land-observation orchestration onboard instead of treating the spacecraft as a passive camera \cite{nasa-challenge-brief}. The target use case is regenerative agriculture, where disease and irrigation stress must be detected early enough to enable patch-level intervention instead of blanket chemical application. Current land products are typically too slow or too coarse for that workflow, even when the underlying algorithms are scientifically strong \cite{gold-webinar1,mod13a1,mod11a1}.

The proposed contribution is a two-stage, MDP-formulated scheduling policy executed as a deterministic, flight-verifiable rule set on a LEON4FT processor. A low-cost spectral screen runs every pass; a more expensive thermal confirmation step runs only when the belief state indicates likely stress. The payload is framed as a hosted module on Loft Orbital's YAM platform and hands alerts to existing agricultural software through OpenET-linked ground services \cite{mytkowicz-webinar2,openet-api}. The rest of the paper focuses on the judge-facing questions of creativity, feasibility, impact, business viability, and next-step readiness.
